\section*{Chapitre \ref{ch:b2:vertebrate-development} --- Le développement embryonnaire des vertébrés}

\begin{solution}{exo:b2:vertebrate-development:1}
\omterm{def:b2:vertebrate-development:egg}{Segmentation}~: mitoses rapides sans croissance, qui produisent la
\omterm{def:b2:vertebrate-development:blastula}{blastula}. \omterm{def:b2:vertebrate-development:blastula}{Blastula}~: sphère creuse ou disque de petites cellules autour
du \omterm{def:b2:vertebrate-development:blastula}{blastocœle}. Gastrulation~: les mouvements cellulaires qui font entrer
le mésoderme et l'endoderme, et produisent la gastrula à trois feuillets,
pourvue d'une cavité intestinale et d'une notocorde. Neurulation~: le
reploiement de l'ectoderme dorsal en \omterm{prop:b2:vertebrate-development:neurulation}{tube neural}, accompagné de la
\omterm{def:b2:vertebrate-development:egg}{segmentation} des somites, qui produit la neurula dotée du plan
d'organisation des vertébrés.
\end{solution}

\begin{solution}{exo:b2:vertebrate-development:2}
Ectoderme~: épiderme, encéphale et moelle épinière, nerfs périphériques
(\omterm{prop:b2:vertebrate-development:neurulation}{crête neurale}), cristallin et cellules pigmentaires. Mésoderme~:
notocorde, muscles, os et cartilage, rein, cœur et sang, derme.
Endoderme~: épithélium intestinal, foie, pancréas, poumons, thyroïde.
\end{solution}

\begin{solution}{exo:b2:vertebrate-development:3}
Peu de \omterm{def:b2:vertebrate-development:egg}{vitellus}~: \omterm{def:b2:vertebrate-development:egg}{segmentation} totale et presque égale, et gastrulation
par involution à travers un blastopore circulaire dans une sphère creuse
(grenouille~: \omterm{def:b2:vertebrate-development:egg}{segmentation} totale mais inégale, les cellules végétatives
chargées de \omterm{def:b2:vertebrate-development:egg}{vitellus} étant grosses et lentes). Beaucoup de \omterm{def:b2:vertebrate-development:egg}{vitellus}
(poulet)~: \omterm{def:b2:vertebrate-development:egg}{segmentation} limitée à un disque posé sur le \omterm{def:b2:vertebrate-development:egg}{vitellus}, et
gastrulation à travers une fente, la ligne primitive, dans le disque
plat, les feuillets s'étalant sur le \omterm{def:b2:vertebrate-development:egg}{vitellus} au lieu d'enclore une
cavité.
\end{solution}

\begin{solution}{exo:b2:vertebrate-development:4}
L'\omterm{def:b2:vertebrate-development:induction}{induction} est la modification du destin d'un tissu compétent par un
signal issu d'un tissu voisin. Les cellules végétatives induisent le
mésoderme dans la zone marginale sus-jacente~; l'\omterm{def:b2:vertebrate-development:induction}{organisateur} (mésoderme
de la lèvre dorsale) induit la plaque neurale dans l'ectoderme dorsal~;
la vésicule optique induit le cristallin dans l'ectoderme céphalique.
\end{solution}

\begin{solution}{exo:b2:vertebrate-development:5}
$2^{12} = 4096$ cellules de rayon $1/2^{4} = \qty{62.5}{\micro m}$~;
surface multipliée par $2^{4} = 16$. La surface supplémentaire est la
membrane des épithéliums que la gastrulation repliera en peau et en
intestin, et la surface par laquelle la \omterm{def:b2:vertebrate-development:blastula}{blastula} échange avec son
milieu.
\end{solution}

\begin{solution}{exo:b2:vertebrate-development:6}
$10^{-4}\times 2^{k} = 0.4$~: $2^{k} = 4000$, $k = 12$. Avec quatre fois
plus d'ADN, le rapport part de $4\times 10^{-4}$ et atteint $0.4$
pour $2^{k} = 1000$~: à la dixième division, deux plus tôt.
\end{solution}

\begin{solution}{exo:b2:vertebrate-development:7}
$k = \ln 2/600 = \qty{1.16e-3}{s^{-1}}$~; $\lambda = \sqrt{2\times
10^{-12}/1.16\times 10^{-3}} = \qty{42}{\micro m}$. Seuils à
$x = \lambda\ln 2 = \qty{29}{\micro m}$ et $\lambda\ln 10 =
\qty{96}{\micro m}$~: des bandes de 29, de 67 et le reste.
\end{solution}

\begin{solution}{exo:b2:vertebrate-development:8}
Doubler la source déplace les deux frontières vers l'extérieur de
$\lambda\ln
2 = \qty{29}{\micro m}$~: la première bande double, la deuxième garde sa
largeur. Doubler $k$ divise $\lambda$ par $\sqrt 2$
(\qty{29}{\micro m}) et, comme $c_0 = j/\sqrt{Dk}$, divise aussi $c_0$
par $\sqrt 2$~: les frontières passent à $\lambda'\ln(c_0'/T) =
29\times(0.69 - 0.35) = \qty{10}{\micro m}$ et $29\times(2.30 -
0.35) = \qty{57}{\micro m}$ --- les deux bandes rétrécissent.
\end{solution}

\begin{solution}{exo:b2:vertebrate-development:9}
Une lèvre dorsale prélevée sur une gastrula de triton non pigmentée fut
greffée sur la face ventrale d'une gastrula pigmentée~; elle s'enroula
vers l'intérieur et l'hôte forma un second embryon --- \omterm{prop:b2:vertebrate-development:neurulation}{tube neural},
somites, intestin --- soudé au premier. La différence de pigmentation
montra que le second axe était fait de cellules de l'hôte~: le greffon
les avait donc induites au lieu de construire lui-même l'axe.
L'expérience exclut l'autodifférenciation du greffon comme explication,
et montra qu'une petite région porte le signal qui organise l'axe tout
entier.
\end{solution}

\begin{solution}{exo:b2:vertebrate-development:10}
(a) De la peau ventrale~: l'ectoderme précoce est compétent mais pas
encore induit, et il suit son nouvel environnement. (b) Du tissu
nerveux~: dans la gastrula âgée, il a été induit et il est déterminé.
(c) Une plaque neurale~: l'ectoderme ventral de gastrula âgée est encore
compétent, et la notocorde l'induit. (d) Aucune structure dorsale~: une
boule de tissu ventral, puisque rien n'induit le mésoderme vers des
destins dorsaux ni l'ectoderme vers des destins nerveux.
\end{solution}

\begin{solution}{exo:b2:vertebrate-development:11}
À \qty{18}{\celsius}~: $90 + 11\times 30 = \qty{420}{min}$, soit
7 heures. À \qty{10}{\celsius}~: \qty{17.5}{h}. La transition se situe à
la douzième division dans les deux cas, car ce qui la déclenche est le
rapport de l'ADN au cytoplasme, qui dépend du nombre de divisions et non
de leur durée~: l'horloge compte des doublements, non des heures.
\end{solution}

\begin{solution}{exo:b2:vertebrate-development:12}
Dans le développement régulateur, les cellules acquièrent leur destin
par \omterm{def:b2:vertebrate-development:induction}{induction} de la part de leurs voisines, si bien que l'embryon peut
se reconstruire après une perte (jumeaux issus d'un œuf séparé,
régénération d'une région greffée)~; dans le développement en mosaïque,
les destins sont hérités avec des déterminants cytoplasmiques localisés,
si bien que chaque blastomère ne forme que sa partie et qu'un embryon
séparé donne deux demi-larves. La différence est de degré~: le croissant
gris de la grenouille est un déterminant, et les embryons en mosaïque
recourent plus tard à des \omterm{def:b2:vertebrate-development:induction}{inductions}~; tout embryon mêle héritage et
conversation, les vertébrés s'appuyant plus longtemps sur la
conversation.
\end{solution}

\begin{solution}{pb:b2:vertebrate-development:1}
\textbf{1.} Œuf $\tfrac43\pi(0.7)^{3} = \qty{1.44}{mm^3}$~; noyau
$\tfrac43\pi(0.03)^{3} = \qty{1.1e-4}{mm^3}$~; rapport $8\times
10^{-5}$.
\textbf{2.} $2^{k} = 4000$~: $k = 12$, 4096 cellules.
\textbf{3.} $75 + 11\times 30 = \qty{405}{min}$, environ 6.75 heures.
\textbf{4.} Rayon $0.7/2^{4} = \qty{44}{\micro m}$~; rapport $8\times
10^{-5}\times 4096 = 0.32$ --- celui d'une cellule ordinaire~: les noyaux
ont rattrapé le cytoplasme.
\textbf{5.} $2^{4} = 16$.
\textbf{6.} Environ $4095\,c$ d'ADN~; l'œuf contient assez d'histones,
de polymérases et de nucléotides en réserve pour douze divisions, si
bien qu'aucune transcription n'est nécessaire.
\textbf{7.} Les gènes zygotiques s'activent, les divisions ralentissent
et perdent leur synchronie, les cellules deviennent mobiles. L'injection
d'ADN supplémentaire avance la transition d'autant de divisions que
l'ADN est en excès, ce qu'un compteur de divisions ou de temps ne
pourrait pas expliquer.
\textbf{8.} À l'opposé du point d'entrée du spermatozoïde, au niveau du
croissant gris fixé par la rotation du cortex de l'œuf après la
fécondation.
\textbf{9.} D'abord le mésoderme dorsal (plaque préchordale, puis
notocorde), qui gagne le plafond de l'archentéron sur la ligne médiane~;
le mésoderme latéral (somites, lame latérale) à côté~; l'endoderme borde
l'archentéron~; l'ectoderme s'étale sur l'extérieur par épibolie.
\textbf{10.} $\qty{2000}{\micro m}/\qty{600}{min} = \qty{3.3}{\micro
m/min}$~: trois fois un fibroblaste --- un feuillet coordonné se déplace
plus vite qu'une cellule isolée.
\textbf{11.} Les mouvements cellulaires exigent des produits de gènes
zygotiques (nouvelles molécules d'adhérence, motilité) et des cycles
ralentis~; avant la transition, les cellules ne peuvent que se diviser.
\textbf{12.} Tout cordé possède une notocorde à un certain stade, et les
vertébrés construisent la colonne vertébrale autour d'elle~; le \omterm{prop:b2:vertebrate-development:neurulation}{tube
neural} est induit par elle, et il est la conséquence du plan, non son
origine.
\textbf{13.} Sans contraction des cellules en bouteille, aucun
blastopore ne se forme, rien ne s'enroule vers l'intérieur et l'embryon
reste une boule dont les feuillets restent en place~: ni intestin, ni
notocorde, ni axe.
\textbf{14.} $k = \ln 2/1200 = \qty{5.8e-4}{s^{-1}}$~; $\lambda =
\sqrt{10^{-12}/5.8\times 10^{-4}} = \qty{42}{\micro m}$.
\textbf{15.} $x_1 = 42\ln 2.5 = \qty{38}{\micro m}$~; $x_2 = 42\ln 12.5
= \qty{105}{\micro m}$.
\textbf{16.} Des bandes de \qty{38}{\micro m} (2.5 cellules),
\qty{67}{\micro m} (4.5 cellules), et le reste.
\textbf{17.} $x_1 = 42\ln 1.25 = \qty{9}{\micro m}$~; $x_2 = 42\ln 6.25
= \qty{77}{\micro m}$~: les deux frontières se déplacent vers l'intérieur
de $\lambda\ln 2 =
\qty{29}{\micro m}$~; la première bande passe de 38 à 9, la bande du
milieu garde sa largeur, la bande externe s'élargit.
\textbf{18.} Déplacée avant la lecture, elle adopte le destin A ---
elle lit la concentration de l'endroit où elle se trouve désormais~;
déplacée après, elle garde le destin B --- elle était déterminée.
\textbf{19.} $x_T = \lambda\ln(c_0/T)$~: une variation de la source d'un
facteur $f$ déplace chaque frontière de $\lambda\ln f$~; une erreur d'un
facteur deux ne décale donc l'organisation que de $0.7\lambda$, une
fraction d'un champ qui s'étend sur plusieurs $\lambda$ --- l'organisation
est presque insensible à la quantité exacte de morphogène.
\textbf{20.} Lèvre dorsale sur la face ventrale~: un second axe. Zone
marginale ventrale sur la face dorsale~: elle est respécifiée par son
environnement et forme du tissu dorsal~; rien de plus.
\textbf{21.} Dans une neurula, l'ectoderme n'est plus compétent~: le
greffon forme une notocorde mais n'induit pas de second \omterm{prop:b2:vertebrate-development:neurulation}{tube neural}.
\textbf{22.} Séparation dans le plan du croissant~: deux embryons
complets~; perpendiculairement~: un embryon et une boule de tissu
ventral.
\textbf{23.} Seule~: une boule épidermique. Avec des cellules
végétatives~: du mésoderme --- muscle, notocorde --- induit dans la
calotte.
\textbf{24.} Pour prouver que le second axe avait été induit dans des
cellules de l'hôte et non construit par le greffon~; si le greffon avait
été indiscernable, l'autodifférenciation n'aurait pas pu être exclue.
Une greffe entre espèces différentes fonctionne (ils utilisèrent deux
espèces de tritons), et c'est ce qui rendait les cellules
distinguables.
\textbf{25.} Douzième division, 4096 cellules, vers 6.75 heures~; des
bandes de 38 et \qty{67}{\micro m}, et tout ce qui se trouve au-delà de
\qty{105}{\micro m}.
\end{solution}
